% Préambule du rapport (pandoc -H preamble.tex) : mise en page moderne et sobre.
% Couleurs reprises de l'application WWI Insight.
\usepackage{xcolor}
\definecolor{accent}{HTML}{0B6E99}
\definecolor{ink}{HTML}{0F1B2D}
\definecolor{muted}{HTML}{5A6678}
\definecolor{codebg}{HTML}{F3F5F8}
\definecolor{hair}{HTML}{DDE2E9}

% Polices
% (fichiers fournis dans fonts/, licence SIL OFL : le rapport compile sans installation de police)
\setmainfont{IBMPlexSans}[Path=fonts/, Extension=.ttf, UprightFont=*-Regular, ItalicFont=*-Italic,
  BoldFont=*-SemiBold, BoldItalicFont=*-SemiBoldItalic]
\setmonofont{IBMPlexMono}[Path=fonts/, Extension=.ttf, Scale=0.88, UprightFont=*-Regular, ItalicFont=*-Italic, BoldFont=*-SemiBold]
\newfontfamily\headfont{IBMPlexSans}[Path=fonts/, Extension=.ttf, UprightFont=*-Regular, ItalicFont=*-Italic, BoldFont=*-SemiBold]

\usepackage{graphicx}
\usepackage{fancyhdr}
\usepackage{titlesec}
\usepackage{titletoc}
\usepackage{caption}
\usepackage{booktabs}
\usepackage{colortbl}
\usepackage{float}
\usepackage{enumitem}

\color{ink}
\linespread{1.18}
\setlength{\parindent}{0pt}
\setlength{\parskip}{0.7em}
\setlist{itemsep=0.25em, topsep=0.3em}
\floatplacement{figure}{H}
\arrayrulecolor{hair}

% En-têtes : discrets, sans filet
\pagestyle{fancy}
\fancyhf{}
\fancyhead[L]{\footnotesize\color{muted}\nouppercase{\leftmark}}
\fancyhead[R]{\footnotesize\color{muted}\thepage}
\renewcommand{\headrulewidth}{0pt}
\fancypagestyle{plain}{\fancyhf{}\fancyfoot[R]{\footnotesize\color{muted}\thepage}}
\renewcommand{\chaptermark}[1]{\markboth{\thechapter.\ #1}{}}

% Titres
\titleformat{\chapter}[display]
  {\headfont\color{ink}}
  {\normalsize\color{accent}\bfseries CHAPITRE \thechapter}
  {6pt}
  {\fontsize{26}{31}\selectfont\bfseries}
  [\vspace{4pt}{\color{hair}\rule{\linewidth}{0.8pt}}]
\titleformat{name=\chapter, numberless}[display]
  {\headfont\color{ink}}{}{0pt}{\fontsize{26}{31}\selectfont\bfseries}
  [\vspace{4pt}{\color{hair}\rule{\linewidth}{0.8pt}}]
\titlespacing*{\chapter}{0pt}{0pt}{28pt}
\titleformat{\section}{\headfont\Large\bfseries\color{ink}}{\color{accent}\thesection}{0.8em}{}
\titleformat{\subsection}{\headfont\large\bfseries\color{ink}}{\color{accent}\thesubsection}{0.8em}{}
\titlespacing*{\section}{0pt}{1.6em}{0.5em}
\titlespacing*{\subsection}{0pt}{1.2em}{0.3em}

% Table des matières
\titlecontents{chapter}[0pt]{\addvspace{0.8em}\bfseries}
  {\color{accent}\thecontentslabel\quad}{}{\hfill\contentspage}
\titlecontents{section}[1.6em]{\addvspace{0.15em}}
  {\color{muted}\thecontentslabel\quad}{}{\titlerule*[0.6pc]{\color{hair}.}\contentspage}

% Légendes : petites, grises, numéro en couleur
\captionsetup{font={small,color=muted}, labelfont={bf,color=accent}, labelsep=quad,
  justification=raggedright, singlelinecheck=false, skip=6pt}

% Code : fond gris clair, sans cadre
\lstset{
  basicstyle=\ttfamily\footnotesize\color{ink},
  backgroundcolor=\color{codebg},
  frame=single, rulecolor=\color{codebg}, framesep=7pt, framerule=0pt,
  xleftmargin=7pt, xrightmargin=7pt,
  numbers=left, numberstyle=\tiny\color{muted}, numbersep=12pt,
  breaklines=true, breakatwhitespace=true, breakindent=1.5em,
  columns=fullflexible, keepspaces=true, showstringspaces=false, tabsize=4,
  keywordstyle=\color{accent}, commentstyle=\itshape\color{muted}, stringstyle=\color{ink},
  captionpos=t, aboveskip=1.1em, belowskip=0.8em,
  literate={é}{{\'e}}1 {è}{{\`e}}1 {ê}{{\^e}}1 {à}{{\`a}}1 {â}{{\^a}}1 {ç}{{\c c}}1
           {ô}{{\^o}}1 {î}{{\^i}}1 {ù}{{\`u}}1 {û}{{\^u}}1 {É}{{\'E}}1 {È}{{\`E}}1
           {’}{{'}}1 {«}{{\guillemotleft}}1 {»}{{\guillemotright}}1 {…}{{...}}1 {·}{{$\cdot$}}1
}
\lstdefinelanguage{yaml}{keywords={true,false,null}, sensitive=false,
  comment=[l]{\#}, morestring=[b]', morestring=[b]"}
\renewcommand{\lstlistingname}{Code}
\renewcommand{\lstlistlistingname}{Liste des codes}

% Noms en français
\AtBeginDocument{%
  \renewcommand{\tablename}{Tableau}%
  \renewcommand{\listtablename}{Liste des tableaux}%
  \hypersetup{colorlinks=true, linkcolor=ink, urlcolor=accent, citecolor=accent}%
}

% Année universitaire affichée sur la page de garde
\newcommand{\ANNEE}{2025–2026}
